\documentclass[11pt]{article}

% Change "review" to "final" to generate the final (sometimes called camera-ready) version.
% Change to "preprint" to generate a non-anonymous version with page numbers.
\usepackage[]{acl}
\usepackage{times}
\usepackage{latexsym}
\usepackage{url}

% For proper rendering and hyphenation of words containing Latin characters (including in bib files)
\usepackage[T1]{fontenc}
% For Vietnamese characters
% \usepackage[T5]{fontenc}
% See https://www.latex-project.org/help/documentation/encguide.pdf for other character sets

% This assumes your files are encoded as UTF8
\usepackage[utf8]{inputenc}
\usepackage{microtype}
\usepackage{inconsolata}
\usepackage{graphicx}

% If the title and author information does not fit in the area allocated, uncomment the following
%
%\setlength\titlebox{<dim>}
%
% and set <dim> to something 5cm or larger.

\title{Group 28 Progress Report:\\Spoiler Review Classifier}


\author{Sky Quan, Second Author, Third Author \\
  \texttt{\{quans6@mcmaster.ca,macid2,macid3\}@mcmaster.ca} }

%\author{
%  \textbf{First Author\textsuperscript{1}},
%  \textbf{Second Author\textsuperscript{1,2}},
%  \textbf{Third T. Author\textsuperscript{1}},
%  \textbf{Fourth Author\textsuperscript{1}},
%\\
%  \textbf{Fifth Author\textsuperscript{1,2}},
%  \textbf{Sixth Author\textsuperscript{1}},
%  \textbf{Seventh Author\textsuperscript{1}},
%  \textbf{Eighth Author \textsuperscript{1,2,3,4}},
%\\
%  \textbf{Ninth Author\textsuperscript{1}},
%  \textbf{Tenth Author\textsuperscript{1}},
%  \textbf{Eleventh E. Author\textsuperscript{1,2,3,4,5}},
%  \textbf{Twelfth Author\textsuperscript{1}},
%\\
%  \textbf{Thirteenth Author\textsuperscript{3}},
%  \textbf{Fourteenth F. Author\textsuperscript{2,4}},
%  \textbf{Fifteenth Author\textsuperscript{1}},
%  \textbf{Sixteenth Author\textsuperscript{1}},
%\\
%  \textbf{Seventeenth S. Author\textsuperscript{4,5}},
%  \textbf{Eighteenth Author\textsuperscript{3,4}},
%  \textbf{Nineteenth N. Author\textsuperscript{2,5}},
%  \textbf{Twentieth Author\textsuperscript{1}}
%\\
%\\
%  \textsuperscript{1}Affiliation 1,
%  \textsuperscript{2}Affiliation 2,
%  \textsuperscript{3}Affiliation 3,
%  \textsuperscript{4}Affiliation 4,
%  \textsuperscript{5}Affiliation 5
%\\
%  \small{
%    \textbf{Correspondence:} \href{mailto:email@domain}{email@domain}
%  }
%}

\begin{document}
\maketitle
% \begin{abstract}
% \end{abstract}

\section{Introduction}

Many people browse reviews before starting a book. There are popular platforms like Goodreads to facilitate this, but reviews often contain major spoilers. This can ruin the experience of reading the book and upset readers that just want to hear a high-level overview or opinions. Goodreads lets reviewers hide this content behind spoiler tags, but it is based on trust, since humans decide whether to include the spoiler tags or not. Because of this, readers can have a book spoiled before they ever open it even if it doesn't have a spoiler tag.  So, creating a spoiler detector keeps readers safe by checking for untagged spoilers before they are posted, letting readers safely filter through reviews that don't have spoilers. This would make review platforms safer for readers to get a high-level overview of a book without worrying about being spoiled.

\section{Related Work}

Here, talk about the related work you encountered for your approach. Cite at least 5 references. Refer to item 2. No one has done exactly your task? Write about the most similar thing you can find. This should be around 0.25-0.5 pages.

\section{Dataset}

The source dataset is the Goodreads Spoilers Dataset, originally collected and released by McAuley Lab at UC San Diego, and reposted on Kaggle by Ahmad Mujtaba.
\begin{itemize}
  \item Original source: 
  \item \url{https://cseweb.ucsd.edu/~jmcauley/datasets/goodreads.html}
  \item Kaggle link: 
  \item \url{https://www.kaggle.com/datasets/pypiahmad/goodreads-book-review}
\end{itemize}
The data was collected in late 2017 from goodreads.com by scraping users' public Goodreads shelves. Given the dataset size of 1,378,033 reviews (of 25,475 books by 18,892 users), we use a smaller, randomly sampled subset of data composed of 10,000.
 
To obtain this smaller and randomized subset, we used stratified reservoir sampling. This was used for two main reasons. First, to give every datapoint a randomized chance to appear in our smaller subset. Secondly, to combat the class imbalance within the source dataset. Since the dataset contains many more reviews without spoilers compared to those with, randomly sampling datapoints doesn't give the model enough spoilers to train on. So, stratified reservoir sampling allows us to also sample a predefined number of spoilers and non-spoilers. With the parameters we set, we create a dataset with 10,000 reviews where 20\% of the reviews contain spoilers and 80\% of the reviews do not, regardless of class imbalance.
 
The source dataset does not label each review as containing or spoiler or not. However, every spoiler within the dataset contains the text “(view spoiler)[\ \ ... (hide spoiler)]”. Where ... represents the specific spoiler to each book. So, if a review contains the text “(view spoiler)” or “(hide spoiler)”, it was labelled as 1, contains a spoiler, and if it did not contain the text, it was labelled as 0, contains no spoiler. The labels were placed under the newly created column has\_spoiler.
 
The dataset is then cleaned by removing the unnecessary columns ['userid', 'rating', 'date\_added', 'date\_updated', 'read\_at', 'started\_at', 'n\_votes', 'n\_comments'] while keeping the necessary columns ['review\_id', 'book\_id', 'review\_text', 'has\_spoiler']. The spoiler tags “(view spoiler)[“ and “(hide spoiler)]” are then removed from review\_text. Finally, the dataset is split into training, validation, and testing data based on an 80:10:10 ratio, respectively.
 
\begin{verbatim}
Example datapoints, review_text shortened for readibility:

{'book_id': '16981',
'review_id': 'a5d2c3628987712d0e05c4f90798eb67', 
'review_text': 'Recommended by Don Katz. Avail for free in December...', 
'has_spoiler': 0} 

{'book_id': '28684704', 
'review_id': '2ede853b14dc4583f96cf5d120af636f', 
'review_text': '... can't be discussed without spoilers ... Daniela says, ...', 
'has_spoiler': 1} 
\end{verbatim}

\section{Features}

Describe any features you used for your model, or how your data was input to your model. Are you doing feature engineering or feature selection? Are you learning embeddings? Is it all part of one neural network? Refer to item 2. This may range from 0.25 pages to 0.5 pages.

\section{Implementation}

Describe your model and implementation here. Refer to item 4. This may take around a page.

\section{Results and Evaluation}

How are you evaluating your model? What results do you have so far? What are your baselines? Refer to item 5. This may take around 0.5 pages.

\section{Feedback and Plans}

Write about your plans for the remainder of the project. This should include a discussion of the feedback you received from your TA, and how you plan to improve your approach. Reflect on your implementation and areas for improvement. Refer to item 6. This may be around 0.5 pages.

\section{Template Notes}

You can remove this section or comment it out, as it only contains instructions for how to use this template. You may use subsections in your document as you find appropriate.

\subsection{Tables and figures}

See Table~\ref{citation-guide} for an example of a table and its caption.
See Figure~\ref{fig:experiments} for an example of a figure and its caption.


\begin{figure}[t]
  \includegraphics[width=\columnwidth]{example-image-golden}
  \caption{A figure with a caption that runs for more than one line.
    Example image is usually available through the \texttt{mwe} package
    without even mentioning it in the preamble.}
  \label{fig:experiments}
\end{figure}

\begin{figure*}[t]
  \includegraphics[width=0.48\linewidth]{example-image-a} \hfill
  \includegraphics[width=0.48\linewidth]{example-image-b}
  \caption {A minimal working example to demonstrate how to place
    two images side-by-side.}
\end{figure*}


\subsection{Citations}

\begin{table*}
  \centering
  \begin{tabular}{lll}
    \hline
    \textbf{Output}           & \textbf{natbib command} & \textbf{ACL only command} \\
    \hline
    \citep{Gusfield:97}       & \verb|\citep|           &                           \\
    \citealp{Gusfield:97}     & \verb|\citealp|         &                           \\
    \citet{Gusfield:97}       & \verb|\citet|           &                           \\
    \citeyearpar{Gusfield:97} & \verb|\citeyearpar|     &                           \\
    \citeposs{Gusfield:97}    &                         & \verb|\citeposs|          \\
    \hline
  \end{tabular}
  \caption{\label{citation-guide}
    Citation commands supported by the style file.
  }
\end{table*}

Table~\ref{citation-guide} shows the syntax supported by the style files.
We encourage you to use the natbib styles.
You can use the command \verb|\citet| (cite in text) to get ``author (year)'' citations, like this citation to a paper by \citet{Gusfield:97}.
You can use the command \verb|\citep| (cite in parentheses) to get ``(author, year)'' citations \citep{Gusfield:97}.
You can use the command \verb|\citealp| (alternative cite without parentheses) to get ``author, year'' citations, which is useful for using citations within parentheses (e.g. \citealp{Gusfield:97}).

\subsection{References}

\nocite{Ando2005,andrew2007scalable,rasooli-tetrault-2015}

Many websites where you can find academic papers also allow you to export a bib file for citation or bib formatted entry. Copy this into the \texttt{custom.bib} and you will be able to cite the paper in the \LaTeX{}. You can remove the example entries.

\subsection{Equations}

An example equation is shown below:
\begin{equation}
  \label{eq:example}
  A = \pi r^2
\end{equation}

Labels for equation numbers, sections, subsections, figures and tables
are all defined with the \verb|\label{label}| command and cross references
to them are made with the \verb|\ref{label}| command.
This an example cross-reference to Equation~\ref{eq:example}. You can also write equations inline, like this: $A=\pi r^2$.


% \section*{Limitations}

\section*{Team Contributions}

Write in this section a few sentences describing the contributions of each team member. What did each member work on? Refer to item 7.

Sky Quan: Handled everything involving the dataset. Loaded the source dataset and completed all the preprocessing: cleaning, labelling, and spliting the dataset as required. Wrote the Introduction and Dataset section.

% Bibliography entries for the entire Anthology, followed by custom entries
%\bibliography{custom,anthology-overleaf-1,anthology-overleaf-2}

% Custom bibliography entries only
\bibliography{custom}

% \appendix

% \section{Example Appendix}
% \label{sec:appendix}

% This is an appendix.

\end{document}
